\chapter*{Abstract}
\addcontentsline{toc}{chapter}{Abstract}

Brzozowski derivatives and Antimirov partial derivatives
enable incremental regular-expression processing and
provide a basis for reconstructing parse trees.
Building on Sulzmann and Lu's work on POSIX parsing
\citep{sulzmannlu2014} and non-POSIX submatching
\citep{sulzmannlu2012}, this thesis investigates the
implementation and performance of derivative-based
regular-expression parsers.

Four parsers are implemented in Rust: two using
Brzozowski derivatives and two using Antimirov
partial derivatives.
Each approach includes a standard variant using
injection functions and an optimized bit-coded variant that
records parsing decisions for later reconstruction.
The derivative variants construct POSIX parse trees,
while the partial-derivative variants construct
full trees under non-POSIX disambiguation.
Execution time, memory usage, and translation cost
are evaluated on a real corpus and a synthetic
stress-test dataset.

The results show that partial-derivative parsers
generally perform best, although rankings vary
across workloads.
Bit-coding improves the standard-derivative parser's
execution time by approximately four to five times
and reduces peak memory by roughly two orders
of magnitude, but offers smaller and less consistent
benefits for partial derivatives.
Overall, performance depends on the derivative
construction, reconstruction representation,
and characteristics of the input.